\documentclass[xetex,mathserif,serif]{beamer}

\usepackage{xunicode}
\usepackage{xltxtra}
\usepackage{color}
\usepackage{url}
\usepackage{listings}
\usepackage{fontspec}
\usepackage{geometry}
\usepackage{lastpage}
\usepackage{fancyhdr}
\usepackage{amsmath}
\usepackage{amsthm}
\usepackage{amssymb}
\usepackage{blkarray}
\usepackage{multicol}
\usepackage{relsize}

\definecolor{solarized@base03}{HTML}{002B36}
\definecolor{solarized@base02}{HTML}{073642}
\definecolor{solarized@base01}{HTML}{586e75}
\definecolor{solarized@base00}{HTML}{657b83}
\definecolor{solarized@base0}{HTML}{839496}
\definecolor{solarized@base1}{HTML}{93a1a1}
\definecolor{solarized@base2}{HTML}{EEE8D5}
\definecolor{solarized@base3}{HTML}{FDF6E3}
\definecolor{solarized@yellow}{HTML}{B58900}
\definecolor{solarized@orange}{HTML}{CB4B16}
\definecolor{solarized@red}{HTML}{DC322F}
\definecolor{solarized@magenta}{HTML}{D33682}
\definecolor{solarized@violet}{HTML}{6C71C4}
\definecolor{solarized@blue}{HTML}{268BD2}
\definecolor{solarized@cyan}{HTML}{2AA198}
\definecolor{solarized@green}{HTML}{859900}
\definecolor{yaleblue}{HTML}{0E4C92}

\setbeamertemplate{navigation symbols}{}
% \setbeamerfont{title}{family=\old}
% \setbeamerfont{author}{family=\tfont}%
% \setbeamerfont{frametitle}{family=\oldA}
% \setbeamerfont{date}{family=\dfont}

\setbeamertemplate{itemize items}{--}
\setbeamercolor*{item}{fg=black}

\defaultfontfeatures{Mapping=tex-text}
\hypersetup{pdfstartview={FitH}}

\newcommand{\old}[1]{\fontspec[Alternate=1,Ligatures={Common}]{Hoefler Text}\fontsize{18pt}{30pt}\selectfont #1}%
\newcommand{\oldA}[1]{\fontspec[Alternate=1,Ligatures={Common, Rare}]{Hoefler Text}\fontsize{12pt}{15pt}\selectfont #1}%
\newcommand{\oldB}[1]{\fontspec[Ligatures={Common}]{Didot}\fontsize{12pt}{15pt}\color{solarized@base02}\selectfont #1}%
\newcommand{\tfont}[1]{\fontspec[Alternate=1,Ligatures={Common}]{Hoefler Text}\fontsize{12pt}{20pt}\selectfont #1}%
\newcommand{\dfont}[1]{\fontspec[Ligatures={Common}]{Didot}\fontsize{12pt}{12pt}\selectfont #1}%

\newcommand{\minimize}{\mathop{\mathrm{minimize}}}
\newcommand{\argmin}{\mathop{\mathrm{arg\,min}}}
\newcommand{\argmax}{\mathop{\mathrm{arg\,max}}}
\newcommand{\st}{\mathop{\mathrm{subject\,\,to}}}

\newcommand\independent{\protect\mathpalette{\protect\independenT}{\perp}}
\def\independenT#1#2{\mathrel{\rlap{$#1#2$}\mkern2mu{#1#2}}}

\setlength{\parindent}{0pt}
\setlength{\parskip}{12pt}

\setromanfont [Ligatures={Common}, Numbers={OldStyle}, Variant=01,
 BoldFont={LinLibertine_RB.otf},
 ItalicFont={LinLibertine_RI.otf},
 BoldItalicFont={LinLibertine_RBI.otf}
 ]{LinLibertine_R.otf}



\begin{document}

%%%%%%%%%%%%%%%%%%%%%%%%%%%%%%%%%%%%%%%%%%%%%%%%%%%
\begin{frame}[fragile] \frametitle{}

\vfill

{\fontsize{0.7cm}{0cm}\selectfont Lecture 16 \\\vspace{0.2cm}
Solving GLMs via IRWLS}\\\vspace{0.5cm}
09 November 2015

\vspace{2cm}

\begin{minipage}{0.6\textwidth}
Taylor B. Arnold \\
Yale Statistics \\
STAT 312/612
\end{minipage}
\hfill
\begin{minipage}{0.3\textwidth}\raggedleft
\includegraphics[scale=0.3]{../yale-logo.png}
\end{minipage}%

\end{frame}

%%%%%%%%%%%%%%%%%%%%%%%%%%%%%%%%%%%%%%%%%%%%%%%%%%%
\begin{frame}[fragile] \frametitle{}

{\color{yaleblue}\fontsize{16pt}{20pt}\selectfont Notes}

\begin{itemize}
\item problem set 5 posted; due next class \pause
\item problem set 6, November 18th
\end{itemize}

\end{frame}

%%%%%%%%%%%%%%%%%%%%%%%%%%%%%%%%%%%%%%%%%%%%%%%%%%%
\begin{frame}[fragile] \frametitle{}

{\color{yaleblue}\fontsize{16pt}{20pt}\selectfont Goals for today}

\begin{itemize}
\item fixed PCA example from last time
\item how to solve logistic regression via weighted least squares
\item classification problem on image corpus
\end{itemize}

\end{frame}

%%%%%%%%%%%%%%%%%%%%%%%%%%%%%%%%%%%%%%%%%%%%%%%%%%%
\begin{frame}[fragile] \frametitle{}

{\color{yaleblue}\fontsize{16pt}{20pt}\selectfont fixed PCA example from last time}

\end{frame}

%%%%%%%%%%%%%%%%%%%%%%%%%%%%%%%%%%%%%%%%%%%%%%%%%%%
\begin{frame}[fragile] \frametitle{}

\textbf{GLMs}

Recall that we define generalized linear models such that
the mean of $y$ is some function of $X\beta$, rather than
directly equal to it:
\begin{align*}
\mathbb{E} (y | X) &= g^{-1} \left( X \beta \right)
\end{align*}
With $g$, called the \textit{link function}, equal to some
fixed and known function.

\end{frame}

%%%%%%%%%%%%%%%%%%%%%%%%%%%%%%%%%%%%%%%%%%%%%%%%%%%
\begin{frame}[fragile] \frametitle{}

\textbf{GLMs}

Recall that we define generalized linear models such that
the mean of $y$ is some function of $X\beta$, rather than
directly equal to it:
\begin{align*}
\mathbb{E} (y | X) &= g^{-1} \left( X \beta \right)
\end{align*}
With $g$, called the \textit{link function}, equal to some
fixed and known function.

\end{frame}

%%%%%%%%%%%%%%%%%%%%%%%%%%%%%%%%%%%%%%%%%%%%%%%%%%%
\begin{frame}[fragile] \frametitle{}

\textbf{Logistic regression}

The logistic regression function uses $g$ equal to the logit
function to describe a distribution with $y \in \{0,1\}$.
Specifically, we have the following description of the statistical
model:
\begin{align*}
\mathbb{E} (y | X) &= \text{logit}^{-1} \left( X \beta \right) \\
&= \frac{1}{1 + e^{- X \beta}}
\end{align*}

\end{frame}

%%%%%%%%%%%%%%%%%%%%%%%%%%%%%%%%%%%%%%%%%%%%%%%%%%%
\begin{frame}[fragile] \frametitle{}

Now that we have discussed how to solve the ordinary least squares
equation, you may wonder how we would go about solving the logistic
regression problem.

\end{frame}

%%%%%%%%%%%%%%%%%%%%%%%%%%%%%%%%%%%%%%%%%%%%%%%%%%%
\begin{frame}[fragile] \frametitle{}

Recall that the likelihood of the data point $y_i$ given its
mean $p_i$ is equal to:
\begin{align*}
L_i(y_i | p_i)
&= \text{exp} \left\{y_i \cdot \log\left(\frac{p_i}{1-p_i}\right) + \log(1 - p_i) \right\} \\
&= \text{exp} \left\{y_i \cdot \theta_i + \frac{e^{\theta_i}}{e^{\theta_i} + 1}\right\} \\
&= \text{exp} \left\{y_i \cdot \theta_i - b(\theta_i) \right\}
\end{align*}
Where:
\begin{align*}
\theta_i &= \log\left(\frac{p_i}{1-p_i}\right) \\
b(\theta_i) &= -1 \cdot \frac{e^{\theta_i}}{e^{\theta_i} + 1}
\end{align*}

\end{frame}

%%%%%%%%%%%%%%%%%%%%%%%%%%%%%%%%%%%%%%%%%%%%%%%%%%%
\begin{frame}[fragile] \frametitle{}

The derivative of the log-likelihood is then given by:
\begin{align*}
\frac{\partial}{\partial \theta_i}  l_i(y_i | \theta_i)
&= y_i - b'(\theta_i)
\end{align*}
And
\begin{align*}
\frac{\partial \theta_i}{\partial p_i} &=
  \frac{1-p_i}{p_i} \cdot \frac{1}{1}
\end{align*}

\end{frame}

%%%%%%%%%%%%%%%%%%%%%%%%%%%%%%%%%%%%%%%%%%%%%%%%%%%
\begin{frame}[fragile] \frametitle{}

Recall that the log-likelihood of the data point $y_i$ given its
mean $p_i$ is equal to:
\begin{align*}
l_i(y_i | p_i)
&= \log(1-p_i) + y_i \cdot \text{log} \left( \frac{p_i}{1 - p_i} \right)
\end{align*}
\pause Taking the derivative with respect to $p_i$ yields:
\begin{align*}
\frac{\partial}{\partial p_i} l_i(y_i | p_i)
&= \frac{1}{p_i - 1} + y_i \cdot \frac{1 - p_i}{p_i} \cdot \frac{}{}
\end{align*}

\end{frame}

\end{document}











